\section{Per-Benchmark Detailed Results}
\label{app:per-benchmark}

Table~\ref{tab:per-benchmark} will report the paired delta ($\Delta$, in pp) and negative interference rate (NIR, \%) for each condition on each individual benchmark, complementing the per-benchmark results in Table~\ref{tab:pinchbench-results} and Table~\ref{tab:agentbench-results}.

\textbf{Note:} This table will be populated upon completion of the systematic repair evaluation (T1--T5 recipes and SFT conditions). The planned experimental design is described in \S\ref{sec:harness:conditions}. The current paper only reports Base condition results on PinchBench (Table~\ref{tab:pinchbench-results}) and AgentBench-OpenClaw (Table~\ref{tab:agentbench-results}).

\begin{table}[h]
\centering
\caption{Per-benchmark test results: paired delta ($\Delta$, pp) and NIR (\%). Results will be reported separately for AgentBench-OC and PinchBench upon completion of the repair evaluation pipeline. [Placeholder --- to be filled]}
\label{tab:per-benchmark}
\resizebox{\linewidth}{!}{%
\begin{tabular}{lcccc}
\toprule
 & \multicolumn{2}{c}{AB-OC} & \multicolumn{2}{c}{PB} \\
\cmidrule(lr){2-3}\cmidrule(lr){4-5}
\textbf{Condition} & $\Delta$ & NIR & $\Delta$ & NIR \\
\midrule
Base & 0.0 & 0.0\% & 0.0 & 0.0\% \\
+Memory & \emph{[TBD]} & \emph{[TBD]} & \emph{[TBD]} & \emph{[TBD]} \\
+Control & \emph{[TBD]} & \emph{[TBD]} & \emph{[TBD]} & \emph{[TBD]} \\
+Collab & \emph{[TBD]} & \emph{[TBD]} & \emph{[TBD]} & \emph{[TBD]} \\
+ProcSupport & \emph{[TBD]} & \emph{[TBD]} & \emph{[TBD]} & \emph{[TBD]} \\
+Memory+Control (T5) & \emph{[TBD]} & \emph{[TBD]} & \emph{[TBD]} & \emph{[TBD]} \\
SFT-100 & \emph{[TBD]} & \emph{[TBD]} & \emph{[TBD]} & \emph{[TBD]} \\
SFT-300 & \emph{[TBD]} & \emph{[TBD]} & \emph{[TBD]} & \emph{[TBD]} \\
SFT-1000 & \emph{[TBD]} & \emph{[TBD]} & \emph{[TBD]} & \emph{[TBD]} \\
\bottomrule
\end{tabular}}
\end{table}

\section{Ablation Study for Combined Recipe (T5)}
\label{app:ablation}

Table~\ref{tab:t5-ablation} will report CRR for each individual component (Memory alone, Control alone) and the combined T5 (+Memory+Control), on each failure category. This ablation is planned to be conducted on the test set.

\begin{table}[h]
\centering
\caption{Ablation of T5 (+Memory+Control) against individual components. CRR will be reported for each failure category upon completion of the evaluation. [Placeholder --- to be filled]}
\label{tab:t5-ablation}
\resizebox{0.8\linewidth}{!}{%
\begin{tabular}{lccccc}
\toprule
\textbf{Condition} & \textbf{A} & \textbf{B} & \textbf{C} & \textbf{D} & \textbf{E} \\
\midrule
+Memory & \emph{[TBD]} & \emph{[TBD]} & \emph{[TBD]} & \emph{[TBD]} & \emph{[TBD]} \\
+Control & \emph{[TBD]} & \emph{[TBD]} & \emph{[TBD]} & \emph{[TBD]} & \emph{[TBD]} \\
+Memory+Control (T5) & \emph{[TBD]} & \emph{[TBD]} & \emph{[TBD]} & \emph{[TBD]} & \emph{[TBD]} \\
\bottomrule
\end{tabular}}
\end{table}

\section{SFT Data Construction Details}
\label{app:sft-data}

For the SFT training data construction, corrected trajectories will be generated by prompting GPT-4o with the following template:

\medskip
\fbox{
\begin{minipage}{0.95\linewidth}
\textbf{System:} You are an expert agent that always produces correct trajectories. Given the task below, generate the correct sequence of actions and tool calls that leads to the gold answer. \\
\textbf{User:} Task: [task description] \\
Gold Answer: [gold answer] \\
Agent Trajectory (failed): [original trajectory] \\
Please generate the corrected trajectory with explanations.
\end{minipage}
}
\medskip

\textbf{Verification:} SFT training is planned as the next phase. Trajectory quality verification protocol (including human spot-checks of corrected trajectories) is planned as part of the SFT data construction pipeline.

\section{Annotation Protocol Details}
\label{app:annotation}

Formal failure annotation with inter-annotator agreement measurement is planned as part of the complete evaluation pipeline. The planned protocol follows the two-phase design in \S\ref{sec:taxonomy:annotation}: Phase 1 (pilot) with 60 failure traces and Phase 2 (formal) with 140 additional cases, targeting Cohen's $\kappa \ge 0.70$.

\section{Additional Baseline Details}
\label{app:baselines}

\textbf{Human ceiling:} Human expert evaluation on PinchBench and AgentBench-OC is planned but not yet conducted. This will provide an upper-bound reference for ceiling effect analysis.

\textbf{Inter-condition correlation:} Pearson correlation of paired deltas across conditions is planned as part of the full evaluation. This analysis will confirm that conditions targeting different failure categories capture distinct repair mechanisms.
